
\begin{proposition}[Two commuting semicomplete fields on the actual component cover]\label{cdp:st:C.2.UN.16}
Let $f:U\to\Delta$ be a proper holomorphic map from a smooth
complex threefold, with connected fibres and with smooth complex
two-tori over $\Delta^*$. Fix an irreducible central component $S_i$
of multiplicity $m_i$. Make the actual base change $t=s^{m_i}$,
normalize, and resolve as in \ref{cdp:st:C.2.UN.15}. Every chosen smooth strict
component $E$ above $S_i$ has multiplicity one and carries a
two-dimensional complex vector space
\[
 \mathfrak a_E\subset
 H^0(E,T_E\otimes_{\mathcal O_E}\mathcal M_E)
 \tag{C.CF.1}\label{cdp:eq:C.CF.1}
\]
of commuting meromorphic vector fields. Every member of
$\mathfrak a_E$ is semicomplete on its holomorphic locus.
The map $E\to S_i$ factors through normalization and is
proper and generically finite.

There are two exhaustive geometric alternatives for this particular
space of fields. Its generic rank is two, in which case any complex
basis is a commuting meromorphic frame; or its generic rank is one,
in which case an actual resolution of a meromorphic first integral
and its Stein factor give
\[
 \widehat E\xrightarrow{\ q\ }B,
 \qquad \widehat E\xrightarrow{\rho}E,
 \tag{C.CF.2}\label{cdp:eq:C.CF.2}
\]
where $\rho$ is proper birational, $B$ is a smooth compact curve,
and the connected general fibre of $q$ is rational or elliptic.

When $a(S_i^\nu)=0$, the unique minimal model of $E$ is a
two-torus or a Hopf surface. This last conclusion concerns the
specified generically finite covering component; the maps giving
it are part of the assertion.
\end{proposition}

\begin{proof}
Write $f':U'\to\Delta_s$ for the resolved base change.
The local equations
$w^{m_i}=u$, $z=s/w$ in (C.SC.2), starting from
$t=u z^{m_i}$, prove both multiplicity one and generic
finiteness. Properness and the universal property of normalization
give the asserted factorization $E\to S_i^\nu\to S_i$.
These are the actual maps used below.

Put $\mathcal O=\mathbb C\{s\}$ and
$K=\operatorname{Frac}(\mathcal O)$. Let
\[
 M=(f'_*\mathcal T_{U'/\Delta_s})_0,
 \qquad
 L=\{v\in M\otimes_{\mathcal O}K:
                     \operatorname{ord}_E(v)\ge0\}.
 \tag{C.CF.3}\label{cdp:eq:C.CF.3}
\]
The order in this formula is the minimum of the coefficient
orders in a local tangent frame at the generic point of $E$.
It is independent of that frame. The direct image is coherent
by properness, is torsion-free because its sections are actual
vector fields on $U'$, and has rank two by base change on the
smooth torus locus. Consequently $M$ is a free rank-two
$\mathcal O$-module. Its sections are identified with actual
holomorphic vertical fields by the definition of direct image.
Any two of them commute: their restrictions to a smooth torus
are translation fields, so their bracket vanishes on the dense
punctured tube and hence everywhere.

Here is a finite bound for the saturation in \eqref{cdp:eq:C.CF.3}.
Choose an $\mathcal O$-basis $v_1,v_2$ of $M$.
At a generic point of $E$, multiplicity one makes $s$ a
coordinate, so the two vertical tangent coordinates $x,y$
give
\[
 (v_1\ v_2)=(\partial_x\ \partial_y)A,
 \qquad d=\operatorname{ord}_E(\det A)<\infty.
 \tag{C.CF.4}\label{cdp:eq:C.CF.4}
\]
All entries of $A$ are holomorphic there. The determinant is
nonzero because the evaluation of the two translation fields
is an isomorphism on each nearby torus. In particular $d\ge0$.
For $v=a_1(s)v_1+a_2(s)v_2\in L$, the column $Aa$ has
nonnegative order along $E$. The identity
$a=(\det A)^{-1}\operatorname{adj}(A)(Aa)$ shows that
each $a_j$ has order at least $-d$. The order of a function
of $s$ along $E$ equals its $s$-adic order, again by
multiplicity one. We obtain actual inclusions
\[
                        M\subset L\subset s^{-d}M.
 \tag{C.CF.5}\label{cdp:eq:C.CF.5}
\]
The set $L$ is an $\mathcal O$-submodule. Since $\mathcal O$
is noetherian, \eqref{cdp:eq:C.CF.5} makes it finite; since it is torsion-free
of rank two over the discrete valuation ring $\mathcal O$,
it is free of rank two.

Every member of $L$ is a global meromorphic vertical field
on a sufficiently small representative of $U'$, regular at
the generic point of $E$. Its restriction is therefore a
global meromorphic section of $T_E$. Tangency follows from
$v(s)=0$ at the generic smooth point and then as an identity
of meromorphic sections. Restriction defines the exact map
\[
 \operatorname{res}_E:L\longrightarrow
 H^0(E,T_E\otimes\mathcal M_E),
 \qquad
 \ker(\operatorname{res}_E)=sL.
 \tag{C.CF.6}\label{cdp:eq:C.CF.6}
\]
Indeed, zero restriction means order at least one along $E$;
division by $s$ then belongs to $L$. Conversely, every member
of $sL$ restricts to zero. Thus
\[
                  L/sL\simeq\mathbb C^2
                    \hookrightarrow
                  H^0(E,T_E\otimes\mathcal M_E).
 \tag{C.CF.7}\label{cdp:eq:C.CF.7}
\]
Define $\mathfrak a_E$ to be its image. Since vertical fields
annihilate functions of $s$, brackets of members of
$M\otimes K$ vanish. Hence the fields in \eqref{cdp:eq:C.CF.7} commute.

For completeness, semicompleteness applies to the entire
space, not only to a selected nonzero element. Given
$V\in\mathfrak a_E$, choose its representative $v\in L$.
Its poles are confined to the central fibre. On every
punctured fibre it is a holomorphic translation field
multiplied by constants, and is complete. Remove the
ambient polar divisor. The closedness of the
uniformizability locus for a holomorphic vector field
\cite[Theorem~3.3]{Guillot2024} makes $v$ semicomplete
there, since the punctured tube is dense. The component
$E$ is invariant: locally $v(s)=0$ and factorization of
the local equation of $s$ show that each irreducible
component ideal is invariant wherever $v$ is holomorphic.
Solutions in $E$ inherit single-valuedness. Some removed
curves may be removable poles for $V$; the same closedness
theorem on $E$ minus its actual polar divisor extends
semicompleteness across those curves. This proves the
assertion for every constant linear combination in
$\mathfrak a_E$.

Take a complex basis $V_1,V_2$ of $\mathfrak a_E$.
When $V_1\wedge V_2$ is nonzero, it is the announced
meromorphic frame. Otherwise there is a globally defined
meromorphic function $h$ with
\[
                  V_2=hV_1,
                  \qquad V_1(h)=0.
 \tag{C.CF.8}\label{cdp:eq:C.CF.8}
\]
The local coefficient ratios agree on their overlaps,
which proves that $h$ is global. It is nonconstant by
the complex-linear independence in \eqref{cdp:eq:C.CF.7}. The second
identity follows from the vanishing bracket.
Resolve the graph of $h:E\dashrightarrow\mathbb P^1$
and take the Stein factor of its projection to
$\mathbb P^1$. This gives exactly \eqref{cdp:eq:C.CF.2}, with a
finite map $B\to\mathbb P^1$. The lifted field is
tangent to $q$ and nonzero on a general fibre.
Its semicompleteness passes through the birational
open isomorphism and then over its full holomorphic
locus by the same closedness theorem.

A nonzero semicomplete meromorphic field on a smooth
compact curve has no pole. At a pole of order $r\ge1$
its expression is $z^{-r}a(z)\partial_z$, $a(0)\ne0$.
The time form $z^r a(z)^{-1}dz$ has a primitive
vanishing to order $r+1$; its inverse solution has
distinct analytic continuations around that time
value. They all lie in the punctured coordinate
disc, contradicting single-valuedness. The restricted
field is therefore holomorphic, so
$\deg T_G=2-2g(G)\ge0$ for a general fibre $G$.
This proves the rational-or-elliptic assertion.

Finally suppose $a(S_i^\nu)=0$. The map
$E\to S_i^\nu$ has a finite Stein factor
$E\to T'\to S_i^\nu$, whose first map is
bimeromorphic. For a meromorphic function on $T'$, the
elementary symmetric functions of its values over
the finite map are meromorphic on $S_i^\nu$.
They give a monic algebraic equation over
$\mathcal M(S_i^\nu)$. Thus the meromorphic function
field extension is algebraic and
\[
                            a(E)=a(S_i^\nu)=0.
 \tag{C.CF.9}\label{cdp:eq:C.CF.9}
\]
In particular \eqref{cdp:eq:C.CF.8} cannot occur.

Apply the compact-surface classification of semicomplete
meromorphic fields \cite[Theorem~9.2]{Guillot2024}
separately to $V_1$ and $V_2$. Its fibration alternative
would give a nonconstant meromorphic function on $E$;
therefore each field is holomorphic on some smooth
bimeromorphic model. An algebraic-dimension-zero
surface is not ruled. Its smooth minimal model $E_0$
is unique, and every smooth bimeromorphic model maps
to it by successive point blowdowns. Holomorphic
fields descend under these actual maps: for a
blowdown $\pi$ the equality
$\pi_*\mathcal O=\mathcal O$ sends the derivation
$g\mapsto V(\pi^*g)$ to a holomorphic derivation on
the base. On the common open isomorphism it is the
original field. The two resulting holomorphic fields
on the same $E_0$ remain complex-linearly independent
and commute.

The complete list of minimal compact surfaces
admitting a nonzero holomorphic field is
\cite[Theorem~0.3]{DOT2001}. In algebraic dimension
zero, its possibilities are two-tori, Hopf surfaces,
Inoue surfaces, and minimal class-VII surfaces with
positive second Betti number. The last class has
at most one independent holomorphic field
\cite[pp.~256 and~260]{DOT2000}; Inoue surfaces
also have at most one
\cite[Propositions~2(ii), 3(ii) and~5]{Inoue1974}.
The two independent fields on $E_0$ exclude both
of those alternatives. Hence $E_0$ is a two-torus
or a Hopf surface, as claimed.
\end{proof}


